% Options for packages loaded elsewhere
\PassOptionsToPackage{unicode}{hyperref}
\PassOptionsToPackage{hyphens}{url}
\PassOptionsToPackage{dvipsnames,svgnames,x11names}{xcolor}
%
\documentclass[
  letterpaper,
  DIV=11,
  numbers=noendperiod]{scrartcl}

\usepackage{amsmath,amssymb}
\usepackage{iftex}
\ifPDFTeX
  \usepackage[T1]{fontenc}
  \usepackage[utf8]{inputenc}
  \usepackage{textcomp} % provide euro and other symbols
\else % if luatex or xetex
  \usepackage{unicode-math}
  \defaultfontfeatures{Scale=MatchLowercase}
  \defaultfontfeatures[\rmfamily]{Ligatures=TeX,Scale=1}
\fi
\usepackage{lmodern}
\ifPDFTeX\else  
    % xetex/luatex font selection
\fi
% Use upquote if available, for straight quotes in verbatim environments
\IfFileExists{upquote.sty}{\usepackage{upquote}}{}
\IfFileExists{microtype.sty}{% use microtype if available
  \usepackage[]{microtype}
  \UseMicrotypeSet[protrusion]{basicmath} % disable protrusion for tt fonts
}{}
\makeatletter
\@ifundefined{KOMAClassName}{% if non-KOMA class
  \IfFileExists{parskip.sty}{%
    \usepackage{parskip}
  }{% else
    \setlength{\parindent}{0pt}
    \setlength{\parskip}{6pt plus 2pt minus 1pt}}
}{% if KOMA class
  \KOMAoptions{parskip=half}}
\makeatother
\usepackage{xcolor}
\usepackage[margin=1in]{geometry}
\setlength{\emergencystretch}{3em} % prevent overfull lines
\setcounter{secnumdepth}{5}
% Make \paragraph and \subparagraph free-standing
\makeatletter
\ifx\paragraph\undefined\else
  \let\oldparagraph\paragraph
  \renewcommand{\paragraph}{
    \@ifstar
      \xxxParagraphStar
      \xxxParagraphNoStar
  }
  \newcommand{\xxxParagraphStar}[1]{\oldparagraph*{#1}\mbox{}}
  \newcommand{\xxxParagraphNoStar}[1]{\oldparagraph{#1}\mbox{}}
\fi
\ifx\subparagraph\undefined\else
  \let\oldsubparagraph\subparagraph
  \renewcommand{\subparagraph}{
    \@ifstar
      \xxxSubParagraphStar
      \xxxSubParagraphNoStar
  }
  \newcommand{\xxxSubParagraphStar}[1]{\oldsubparagraph*{#1}\mbox{}}
  \newcommand{\xxxSubParagraphNoStar}[1]{\oldsubparagraph{#1}\mbox{}}
\fi
\makeatother


\providecommand{\tightlist}{%
  \setlength{\itemsep}{0pt}\setlength{\parskip}{0pt}}\usepackage{longtable,booktabs,array}
\usepackage{calc} % for calculating minipage widths
% Correct order of tables after \paragraph or \subparagraph
\usepackage{etoolbox}
\makeatletter
\patchcmd\longtable{\par}{\if@noskipsec\mbox{}\fi\par}{}{}
\makeatother
% Allow footnotes in longtable head/foot
\IfFileExists{footnotehyper.sty}{\usepackage{footnotehyper}}{\usepackage{footnote}}
\makesavenoteenv{longtable}
\usepackage{graphicx}
\makeatletter
\newsavebox\pandoc@box
\newcommand*\pandocbounded[1]{% scales image to fit in text height/width
  \sbox\pandoc@box{#1}%
  \Gscale@div\@tempa{\textheight}{\dimexpr\ht\pandoc@box+\dp\pandoc@box\relax}%
  \Gscale@div\@tempb{\linewidth}{\wd\pandoc@box}%
  \ifdim\@tempb\p@<\@tempa\p@\let\@tempa\@tempb\fi% select the smaller of both
  \ifdim\@tempa\p@<\p@\scalebox{\@tempa}{\usebox\pandoc@box}%
  \else\usebox{\pandoc@box}%
  \fi%
}
% Set default figure placement to htbp
\def\fps@figure{htbp}
\makeatother

\KOMAoption{captions}{tableheading}
\makeatletter
\@ifpackageloaded{caption}{}{\usepackage{caption}}
\AtBeginDocument{%
\ifdefined\contentsname
  \renewcommand*\contentsname{Table of contents}
\else
  \newcommand\contentsname{Table of contents}
\fi
\ifdefined\listfigurename
  \renewcommand*\listfigurename{List of Figures}
\else
  \newcommand\listfigurename{List of Figures}
\fi
\ifdefined\listtablename
  \renewcommand*\listtablename{List of Tables}
\else
  \newcommand\listtablename{List of Tables}
\fi
\ifdefined\figurename
  \renewcommand*\figurename{Figure}
\else
  \newcommand\figurename{Figure}
\fi
\ifdefined\tablename
  \renewcommand*\tablename{Table}
\else
  \newcommand\tablename{Table}
\fi
}
\@ifpackageloaded{float}{}{\usepackage{float}}
\floatstyle{ruled}
\@ifundefined{c@chapter}{\newfloat{codelisting}{h}{lop}}{\newfloat{codelisting}{h}{lop}[chapter]}
\floatname{codelisting}{Listing}
\newcommand*\listoflistings{\listof{codelisting}{List of Listings}}
\makeatother
\makeatletter
\makeatother
\makeatletter
\@ifpackageloaded{caption}{}{\usepackage{caption}}
\@ifpackageloaded{subcaption}{}{\usepackage{subcaption}}
\makeatother

\usepackage{bookmark}

\IfFileExists{xurl.sty}{\usepackage{xurl}}{} % add URL line breaks if available
\urlstyle{same} % disable monospaced font for URLs
\hypersetup{
  pdftitle={Module 2: Philosophies of Education},
  pdfauthor={Nathan Alexander, PhD},
  colorlinks=true,
  linkcolor={blue},
  filecolor={Maroon},
  citecolor={Blue},
  urlcolor={Blue},
  pdfcreator={LaTeX via pandoc}}


\title{Module 2: Philosophies of Education}
\usepackage{etoolbox}
\makeatletter
\providecommand{\subtitle}[1]{% add subtitle to \maketitle
  \apptocmd{\@title}{\par {\large #1 \par}}{}{}
}
\makeatother
\subtitle{EDUC 219 - Howard University}
\author{Nathan Alexander, PhD}
\date{}

\begin{document}
\maketitle

\renewcommand*\contentsname{Table of contents}
{
\hypersetup{linkcolor=}
\setcounter{tocdepth}{3}
\tableofcontents
}

\section{Module overview}\label{module-overview}

This module introduces historical traditions and contemporary discourses
in philosophies of education. We will examine how philosophical ideas
shape beliefs about the purposes of education, the role of schools in
democratic society, the relationships among knowledge, learning,
teaching, and citizenship, and the distribution of educational
opportunity.

\section{Guiding questions}\label{guiding-questions}

\begin{itemize}
\tightlist
\item
  What is a philosophy of education?
\item
  What purposes should education and schools serve in a democratic
  society?
\item
  How have major philosophical traditions shaped American educational
  thought and practice?
\item
  How do different philosophies understand knowledge, learning,
  authority, curriculum, and citizenship?
\end{itemize}

\section{Key concepts}\label{key-concepts}

\begin{itemize}
\tightlist
\item
  Philosophy of education
\item
  Democratic education
\item
  Educational aims and purposes
\item
  Knowledge and curriculum
\item
  Authority and freedom
\item
  Equality, equity, and access
\item
  Diversity and multiculturalism
\item
  Social justice
\item
  Professional identity and agency
\end{itemize}

\begin{center}\rule{0.5\linewidth}{0.5pt}\end{center}

\section{James Banks}\label{james-banks}

\includegraphics[width=0.38\linewidth,height=\textheight,keepaspectratio]{../img/mod01/james-banks.jpg}

\subsection{Multicultural Education}\label{multicultural-education}

James Banks developed a framework for integrating multicultural content
into curricula. His work emphasizes five dimensions:

\begin{itemize}
\tightlist
\item
  Content integration
\item
  Knowledge construction
\item
  Prejudice reduction
\item
  Equity pedagogy
\item
  An empowering school culture and social structure
\end{itemize}

Multicultural education involves more than adding diverse figures or
holidays to the curriculum; it also requires examining how knowledge is
constructed and how school structures distribute opportunity.

\begin{center}\rule{0.5\linewidth}{0.5pt}\end{center}

\section{Django Paris}\label{django-paris}

\subsection{Culturally Sustaining
Pedagogy}\label{culturally-sustaining-pedagogy}

Django Paris extended culturally relevant pedagogy through an explicit
focus on \textbf{sustaining} linguistic, literate, and cultural
pluralism in schooling.

\begin{itemize}
\tightlist
\item
  Centers students' existing cultural and linguistic resources
\item
  Challenges assimilationist assumptions about achievement
\item
  Links pedagogy to positive social transformation
\item
  Treats pluralism as an educational asset rather than a deficit
\end{itemize}

\textbf{Instructional implication:} Teaching should support students in
succeeding academically without requiring them to abandon their
languages, identities, or community practices.

\begin{center}\rule{0.5\linewidth}{0.5pt}\end{center}

\section{Carter G. Woodson}\label{carter-g.-woodson}

\subsection{Black Education}\label{black-education}

Carter G. Woodson advocated for the centrality of Black history in
education.

\begin{itemize}
\tightlist
\item
  Founded Negro History Week, which later developed into Black History
  Month
\item
  Critiqued schooling that devalued Black history, knowledge, and agency
\item
  Developed a lasting critique of the ``mis-education'' of African
  Americans
\item
  Linked historical knowledge to collective self-determination
\end{itemize}

\textbf{Instructional implication:} Curriculum should critically examine
whose histories, achievements, and intellectual traditions are
represented---and whose are excluded.

\begin{center}\rule{0.5\linewidth}{0.5pt}\end{center}

\section{Bob Moses}\label{bob-moses}

\subsection{Literacy Education}\label{literacy-education}

Bob Moses combined civil-rights activism with literacy and mathematics
education.

\begin{itemize}
\tightlist
\item
  Treated literacy, including mathematical literacy, as a civil-rights
  issue
\item
  Viewed quantitative knowledge as a resource for participation and
  empowerment
\item
  Worked with underserved communities through the Algebra Project
\item
  Connected educational opportunity to political and economic access
\end{itemize}

\textbf{Instructional implication:} Access to rigorous mathematics is
not simply an academic matter; it can be understood as a matter of
citizenship, power, and opportunity.

\begin{center}\rule{0.5\linewidth}{0.5pt}\end{center}

\section{Paulo Freire}\label{paulo-freire}

\subsection{Education for Liberation}\label{education-for-liberation}

Paulo Freire developed critical pedagogy, arguing that education should
be a practice of freedom rather than domination.

\begin{itemize}
\tightlist
\item
  Emphasized dialogue rather than one-way transmission
\item
  Distinguished critical inquiry from ``banking'' models of education
\item
  Advanced critical consciousness as an educational aim
\item
  Connected learning to learners' social realities and collective action
\end{itemize}

\textbf{Instructional implication:} Students should be positioned as
knowledgeable participants who can investigate, name, and act on
conditions shaping their lives.

\begin{center}\rule{0.5\linewidth}{0.5pt}\end{center}

\section{Jo-ann Archibald Q'um Q'um
Xiiem}\label{jo-ann-archibald-qum-qum-xiiem}

\subsection{Indigenous Storywork}\label{indigenous-storywork}

Jo-ann Archibald Q'um Q'um Xiiem developed a framework for using
Indigenous stories in education.

Her seven principles are:

\begin{itemize}
\tightlist
\item
  Respect
\item
  Responsibility
\item
  Reciprocity
\item
  Reverence
\item
  Holism
\item
  Interrelatedness
\item
  Synergy
\end{itemize}

\textbf{Instructional implication:} Stories can be understood as sources
of intellectual, ethical, emotional, spiritual, and relational
learning---not merely as illustrative classroom materials.

\begin{center}\rule{0.5\linewidth}{0.5pt}\end{center}

\section{Bryan McKinley Jones
Brayboy}\label{bryan-mckinley-jones-brayboy}

\subsection{Tribal Critical Race
Theory}\label{tribal-critical-race-theory}

Bryan McKinley Jones Brayboy extended Critical Race Theory to address
the particular political, legal, historical, and educational contexts of
Indigenous peoples.

\begin{itemize}
\tightlist
\item
  Centers tribal sovereignty and self-determination
\item
  Examines colonization as enduring rather than merely historical
\item
  Recognizes storytelling as theory-building
\item
  Affirms Indigenous knowledge systems and community-based
  understandings
\end{itemize}

\textbf{Instructional implication:} Equity work involving Indigenous
communities must attend to sovereignty, place, historical relationships,
and Indigenous intellectual traditions.

\begin{center}\rule{0.5\linewidth}{0.5pt}\end{center}

\section{Albert Bandura}\label{albert-bandura}

\subsection{Observational Learning
Theory}\label{observational-learning-theory}

Albert Bandura argued that behavior can be learned through observing
others.

\begin{itemize}
\tightlist
\item
  Learners attend to models and the consequences models receive
\item
  Observation may influence behavior even without direct reinforcement
\item
  Modeling can shape skills, expectations, and beliefs about one's
  capabilities
\item
  Social environments are important sources of learning
\end{itemize}

\textbf{Instructional implication:} Teachers, peers, media, and
classroom norms all provide models that students may observe and
emulate.

\begin{center}\rule{0.5\linewidth}{0.5pt}\end{center}

\section{Asa Hilliard}\label{asa-hilliard}

\subsection{African-Centered
Education}\label{african-centered-education}

Asa Hilliard emphasized the educational importance of cultural identity,
particularly for African American students.

\begin{itemize}
\tightlist
\item
  Advocated for curriculum that includes African and African American
  history and culture
\item
  Critiqued deficit-oriented accounts of Black learners and communities
\item
  Highlighted intellectual traditions that conventional curricula often
  marginalize
\item
  Connected cultural affirmation with rigorous learning
\end{itemize}

\textbf{Instructional implication:} High expectations and cultural
responsiveness should be mutually reinforcing rather than treated as
competing priorities.

\begin{center}\rule{0.5\linewidth}{0.5pt}\end{center}

\section{Abraham Maslow}\label{abraham-maslow}

\subsection{Hierarchy of Needs}\label{hierarchy-of-needs}

Abraham Maslow proposed that people strive to satisfy different
categories of needs.

From lower to higher, the commonly presented levels are:

\begin{enumerate}
\def\labelenumi{\arabic{enumi}.}
\tightlist
\item
  Physiological needs
\item
  Safety needs
\item
  Love and belonging
\item
  Esteem
\item
  Self-actualization
\end{enumerate}

\textbf{Instructional implication:} Learning conditions are shaped by
students' material security, safety, belonging, dignity, and
opportunities for meaningful growth.

\begin{center}\rule{0.5\linewidth}{0.5pt}\end{center}

\section{Carol Dweck}\label{carol-dweck}

\subsection{Growth Mindset}\label{growth-mindset}

Carol Dweck examines how beliefs about intelligence influence motivation
and achievement.

\begin{longtable}[]{@{}
  >{\raggedright\arraybackslash}p{(\linewidth - 4\tabcolsep) * \real{0.3333}}
  >{\raggedright\arraybackslash}p{(\linewidth - 4\tabcolsep) * \real{0.3333}}
  >{\raggedright\arraybackslash}p{(\linewidth - 4\tabcolsep) * \real{0.3333}}@{}}
\toprule\noalign{}
\begin{minipage}[b]{\linewidth}\raggedright
Orientation
\end{minipage} & \begin{minipage}[b]{\linewidth}\raggedright
Core belief
\end{minipage} & \begin{minipage}[b]{\linewidth}\raggedright
Likely response to challenge
\end{minipage} \\
\midrule\noalign{}
\endhead
\bottomrule\noalign{}
\endlastfoot
Fixed mindset & Ability is largely unchangeable & Avoid challenge or
interpret difficulty as evidence of limited ability \\
Growth mindset & Ability can develop through learning, strategy,
feedback, and effort & Treat challenge as an opportunity to learn \\
\end{longtable}

\textbf{Instructional implication:} Feedback can emphasize strategies,
revision, persistence, and learning processes rather than presenting
achievement as evidence of fixed ability.

\begin{center}\rule{0.5\linewidth}{0.5pt}\end{center}

\section{Erik Erikson}\label{erik-erikson}

\subsection{Socioemotional
Development}\label{socioemotional-development}

Erik Erikson described development as a sequence of psychosocial crises
across the lifespan.

\begin{itemize}
\tightlist
\item
  Early development begins with trust versus mistrust
\item
  Later stages involve tensions such as autonomy versus shame,
  initiative versus guilt, and industry versus inferiority
\item
  Adolescence is associated with identity versus role confusion
\item
  Development continues through adulthood
\end{itemize}

\textbf{Instructional implication:} Classroom experiences can influence
students' senses of competence, belonging, identity, and agency.

\begin{center}\rule{0.5\linewidth}{0.5pt}\end{center}

\section{Leon Festinger}\label{leon-festinger}

\subsection{Cognitive Dissonance}\label{cognitive-dissonance}

Leon Festinger's theory of cognitive dissonance proposes that
inconsistencies between beliefs, values, attitudes, and behaviors can
create psychological discomfort.

\begin{itemize}
\tightlist
\item
  Dissonance may motivate a person to change beliefs, behavior, or
  interpretations
\item
  Learners may revise understandings when existing explanations no
  longer fit evidence or experience
\item
  The idea is often connected to constructivist accounts of conceptual
  change
\end{itemize}

\textbf{Instructional implication:} Carefully designed questions,
evidence, and discrepant events can help students recognize and
reconsider incomplete understandings.

\begin{center}\rule{0.5\linewidth}{0.5pt}\end{center}

\section{Sigmund Freud}\label{sigmund-freud}

\subsection{Levels of Consciousness}\label{levels-of-consciousness}

Sigmund Freud proposed that the mind operates across conscious and
unconscious levels.

His structural model includes:

\begin{itemize}
\tightlist
\item
  \textbf{Id:} Primitive motivations and desires
\item
  \textbf{Ego:} The mediating, reality-oriented component
\item
  \textbf{Superego:} Internalized moral standards and conscience
\end{itemize}

\textbf{Instructional implication:} Although not an instructional
theory, Freudian ideas have influenced ways educators and psychologists
think about motivation, emotion, conflict, and development.

\begin{center}\rule{0.5\linewidth}{0.5pt}\end{center}

\section{Gloria Ladson-Billings}\label{gloria-ladson-billings}

\subsection{Culturally Relevant
Pedagogy}\label{culturally-relevant-pedagogy}

Gloria Ladson-Billings developed culturally relevant pedagogy, which
argues that teaching should build from students' cultural backgrounds
and community knowledge.

Culturally relevant teaching seeks to support:

\begin{itemize}
\tightlist
\item
  Academic success
\item
  Cultural competence
\item
  Critical consciousness
\end{itemize}

\textbf{Instructional implication:} Strong teaching advances rigorous
disciplinary learning while supporting students' cultural identities and
their capacity to analyze social inequities.

\begin{center}\rule{0.5\linewidth}{0.5pt}\end{center}

\section{Howard Gardner}\label{howard-gardner}

\subsection{Multiple Intelligences}\label{multiple-intelligences}

Howard Gardner proposed that individuals possess distinct forms of
intelligence, commonly including:

\begin{itemize}
\tightlist
\item
  Linguistic
\item
  Logical-mathematical
\item
  Spatial
\item
  Bodily-kinesthetic
\item
  Musical
\item
  Intrapersonal
\item
  Interpersonal
\end{itemize}

\textbf{Instructional implication:} Teachers can offer varied ways for
students to encounter, represent, practice, and demonstrate learning
rather than relying on a single mode of instruction or assessment.

\begin{center}\rule{0.5\linewidth}{0.5pt}\end{center}

\section{Ivan Pavlov}\label{ivan-pavlov}

\subsection{Classical Conditioning}\label{classical-conditioning}

Ivan Pavlov's classical conditioning model describes the association of
a new stimulus with an existing stimulus-response pairing.

Classic example:

\begin{itemize}
\tightlist
\item
  Food produces salivation in dogs
\item
  A bell is repeatedly paired with food
\item
  The bell eventually produces salivation even when food is absent
\end{itemize}

\textbf{Instructional implication:} Learners can form associations
between classroom conditions, emotional responses, cues, and academic
activities.

\begin{center}\rule{0.5\linewidth}{0.5pt}\end{center}

\section{Jarvis Givens}\label{jarvis-givens}

\subsection{Fugitive Pedagogy}\label{fugitive-pedagogy}

Jarvis Givens examines a hidden tradition of Black teaching developed
under oppressive educational conditions.

\begin{itemize}
\tightlist
\item
  Documents how Black educators navigated systems designed to constrain
  Black learning
\item
  Centers pedagogical care, intellectual resistance, and racial uplift
\item
  Shows how teaching has served as a site of freedom-making
\item
  Connects historical study to current educational practice
\end{itemize}

\textbf{Instructional implication:} Teachers' practices can carry
intellectual and ethical commitments that exceed official curricula,
policies, and institutional constraints.

\begin{center}\rule{0.5\linewidth}{0.5pt}\end{center}

\section{Jean Piaget}\label{jean-piaget}

\subsection{Developmental
Epistemology}\label{developmental-epistemology}

Jean Piaget theorized that children's thinking develops through
qualitatively different stages.

\begin{longtable}[]{@{}
  >{\raggedright\arraybackslash}p{(\linewidth - 4\tabcolsep) * \real{0.3333}}
  >{\raggedright\arraybackslash}p{(\linewidth - 4\tabcolsep) * \real{0.3333}}
  >{\raggedright\arraybackslash}p{(\linewidth - 4\tabcolsep) * \real{0.3333}}@{}}
\toprule\noalign{}
\begin{minipage}[b]{\linewidth}\raggedright
Approximate age
\end{minipage} & \begin{minipage}[b]{\linewidth}\raggedright
Stage
\end{minipage} & \begin{minipage}[b]{\linewidth}\raggedright
General characterization
\end{minipage} \\
\midrule\noalign{}
\endhead
\bottomrule\noalign{}
\endlastfoot
0--2 years & Sensorimotor & Learning through perception and motor
activity \\
3--7 years & Preoperational & Intuitive and symbolic thinking \\
8--11 years & Concrete operational & Logical reasoning about concrete
situations \\
12--15 years and beyond & Formal operational & Abstract and hypothetical
reasoning \\
\end{longtable}

\textbf{Instructional implication:} Learning tasks should be responsive
to learners' developing ways of reasoning, while avoiding assumptions
that age alone fully determines capability.

\begin{center}\rule{0.5\linewidth}{0.5pt}\end{center}

\section{Jerome Bruner}\label{jerome-bruner}

\subsection{Constructivist Theory}\label{constructivist-theory}

Jerome Bruner argued that learners actively construct knowledge by
relating new ideas to their existing schemas or mental models.

\begin{itemize}
\tightlist
\item
  Learning involves making meaning, not merely receiving information
\item
  Prior knowledge shapes how new information is interpreted
\item
  Discovery, inquiry, and revisiting ideas can support understanding
\item
  Instruction can be organized so concepts return with increasing
  sophistication
\end{itemize}

\textbf{Instructional implication:} Teaching should elicit prior ideas,
create opportunities for inquiry, and help students revise and extend
their existing understandings.

\begin{center}\rule{0.5\linewidth}{0.5pt}\end{center}

\section{John Dewey}\label{john-dewey}

\subsection{Learning by Doing}\label{learning-by-doing}

John Dewey emphasized that learning occurs through experience, inquiry,
and reflection.

\begin{itemize}
\tightlist
\item
  Knowledge grows through engagement with meaningful problems
\item
  Experience becomes educative when learners reflect on it
\item
  Schools can prepare students for democratic participation
\item
  Classroom activity should connect to learners' lives and communities
\end{itemize}

\textbf{Instructional implication:} Students learn deeply when they
investigate authentic questions, test ideas, collaborate, and reflect on
consequences.

\begin{center}\rule{0.5\linewidth}{0.5pt}\end{center}

\section{John Hattie}\label{john-hattie}

\subsection{Visible Learning}\label{visible-learning}

John Hattie synthesizes meta-analytic research to identify influences
associated with student achievement.

\begin{itemize}
\tightlist
\item
  Emphasizes the use of evidence in instructional decision-making
\item
  Draws attention to feedback, teacher clarity, formative assessment,
  and visible learning processes
\item
  Encourages educators to examine the impact of their teaching
\item
  Promotes ongoing assessment of whether students are making progress
\end{itemize}

\textbf{Instructional implication:} Teachers can use multiple sources of
evidence to make learning goals, progress, and next steps visible to
students and to inform instruction.

\begin{center}\rule{0.5\linewidth}{0.5pt}\end{center}

\section{Ken Robinson}\label{ken-robinson}

\subsection{Creativity in Education}\label{creativity-in-education}

Ken Robinson advocated for education that nurtures creativity, divergent
thinking, and diverse forms of human talent.

\begin{itemize}
\tightlist
\item
  Critiqued narrow definitions of achievement
\item
  Questioned overly standardized approaches to schooling
\item
  Called for personalized and varied learning opportunities
\item
  Emphasized the relationship between creativity, identity, and
  engagement
\end{itemize}

\textbf{Instructional implication:} Curriculum and assessment can create
room for multiple pathways, imaginative problem-solving, artistic
expression, and students' developing interests.

\begin{center}\rule{0.5\linewidth}{0.5pt}\end{center}

\section{Lawrence Kohlberg}\label{lawrence-kohlberg}

\subsection{Stages of Moral
Development}\label{stages-of-moral-development}

Lawrence Kohlberg described moral reasoning in three broad levels:

\begin{longtable}[]{@{}
  >{\raggedright\arraybackslash}p{(\linewidth - 2\tabcolsep) * \real{0.5000}}
  >{\raggedright\arraybackslash}p{(\linewidth - 2\tabcolsep) * \real{0.5000}}@{}}
\toprule\noalign{}
\begin{minipage}[b]{\linewidth}\raggedright
Level
\end{minipage} & \begin{minipage}[b]{\linewidth}\raggedright
General orientation
\end{minipage} \\
\midrule\noalign{}
\endhead
\bottomrule\noalign{}
\endlastfoot
Pre-conventional & Moral decisions are oriented toward self-interest and
consequences \\
Conventional & Moral decisions are oriented toward social approval,
rules, and local expectations \\
Post-conventional & Moral decisions are oriented toward broader ethical
principles \\
\end{longtable}

\textbf{Instructional implication:} Ethical discussion can invite
students to explain and examine the reasons underlying judgments, rules,
obligations, and competing values.

\begin{center}\rule{0.5\linewidth}{0.5pt}\end{center}

\section{John Locke}\label{john-locke}

\subsection{Tabula Rasa}\label{tabula-rasa}

John Locke's tabula rasa metaphor characterizes the mind as a ``blank
slate'' shaped by experience.

\begin{itemize}
\tightlist
\item
  Emphasizes the formative influence of environment and education
\item
  Historically influenced later empiricist and behaviorist traditions
\item
  Raises questions about the relative roles of experience, teaching, and
  individual development
\end{itemize}

\textbf{Instructional implication:} Educational environments,
experiences, and instruction matter profoundly---but learners should not
be treated as passive recipients of knowledge.

\begin{center}\rule{0.5\linewidth}{0.5pt}\end{center}

\section{Lev Vygotsky}\label{lev-vygotsky}

\subsection{Social Development Theory and
ZPD}\label{social-development-theory-and-zpd}

Lev Vygotsky argued that social interaction is central to cognitive
development.

The \textbf{Zone of Proximal Development (ZPD)} distinguishes between:

\begin{itemize}
\tightlist
\item
  What a learner can do independently
\item
  What a learner can do with assistance from a more capable peer,
  teacher, or other support
\end{itemize}

The ZPD provides a theoretical basis for \textbf{scaffolding}.

\textbf{Instructional implication:} Teachers can design supports,
collaborative structures, prompts, models, and feedback that help
learners do increasingly complex work independently.

\begin{center}\rule{0.5\linewidth}{0.5pt}\end{center}

\section{Carl Rogers}\label{carl-rogers}

\subsection{Experiential Learning}\label{experiential-learning}

Carl Rogers distinguished between academic knowledge and experiential
knowledge.

\begin{itemize}
\tightlist
\item
  Academic knowledge is often externally organized and assessed
\item
  Experiential learning responds to learners' purposes and real-life
  concerns
\item
  Meaningful learning involves personal relevance and active
  participation
\item
  A common example is learning to drive in order to meet a practical
  need
\end{itemize}

\textbf{Instructional implication:} Learning experiences gain power when
students see their relevance, make choices, pursue authentic tasks, and
reflect on their development.

\begin{center}\rule{0.5\linewidth}{0.5pt}\end{center}

\section{B. F. Skinner}\label{b.-f.-skinner}

\subsection{Operant Conditioning}\label{operant-conditioning}

B. F. Skinner argued that learning can be understood as changes in
behavior shaped by consequences.

\begin{itemize}
\tightlist
\item
  Behaviors followed by reinforcement are more likely to recur
\item
  Consequences can strengthen, weaken, or shape responses over time
\item
  Learners can initiate responses rather than merely react to stimuli
\item
  The theory differs from simple stimulus-response association because
  consequences play a central role
\end{itemize}

\textbf{Instructional implication:} Classroom routines, feedback,
incentives, and consequences can influence participation and behavior;
their use should be considered ethically and in relation to long-term
learning goals.

\begin{center}\rule{0.5\linewidth}{0.5pt}\end{center}

\section{Edward Thorndike}\label{edward-thorndike}

\subsection{Connectionism}\label{connectionism}

Edward Thorndike's connectionism proposed that learners form
associations between stimuli and responses.

\begin{itemize}
\tightlist
\item
  Learning develops through trial and error
\item
  Rewarded responses are more likely to be strengthened
\item
  Practice and consequences can influence the formation of habits
\item
  The theory contributed to behaviorist approaches to learning
\end{itemize}

\textbf{Instructional implication:} Repetition, practice, feedback, and
task design can matter for fluency, although complex understanding
generally requires more than reinforced responses.

\begin{center}\rule{0.5\linewidth}{0.5pt}\end{center}

\section{John B. Watson}\label{john-b.-watson}

\subsection{Behaviorism}\label{behaviorism}

John B. Watson argued that much human learning and behavior is shaped by
experience rather than being genetically predetermined.

\begin{itemize}
\tightlist
\item
  Focused on observable behavior as the proper object of scientific
  study
\item
  Minimized attention to unobservable mental processes
\item
  Helped establish behaviorism as an influential psychological tradition
\item
  Emphasized environmental influences on learning and behavior
\end{itemize}

\textbf{Instructional implication:} Observable actions and environmental
conditions are important educational concerns, but behaviorist accounts
may not fully capture learners' meaning-making, identity, or social
context.

\begin{center}\rule{0.5\linewidth}{0.5pt}\end{center}

\section{Max Wertheimer}\label{max-wertheimer}

\subsection{Gestalt Theory}\label{gestalt-theory}

Max Wertheimer's Gestalt theory argues that some ideas can only be
understood as organized wholes rather than as isolated parts.

\begin{itemize}
\tightlist
\item
  The ``whole'' can be more meaningful than the sum of separate elements
\item
  Pattern recognition and structural insight matter in problem-solving
\item
  Understanding often involves seeing relationships among elements
\item
  The theory challenged purely atomistic accounts of learning
\end{itemize}

\textbf{Instructional implication:} Teachers can support learning by
helping students identify patterns, relationships, systems, and the
larger conceptual structure of a problem.

\begin{center}\rule{0.5\linewidth}{0.5pt}\end{center}




\end{document}
